
\documentclass[10pt]{article}
\usepackage[T1]{fontenc}
\usepackage[utf8]{inputenc}
\usepackage[french]{babel}
\usepackage{lmodern,microtype}
\usepackage[a4paper,margin=2.25cm]{geometry}
\usepackage{amsmath,amssymb,mathtools}
\usepackage[dvipsnames]{xcolor}
\usepackage{booktabs,longtable,tabularx,enumitem}
\usepackage[hidelinks]{hyperref}
\definecolor{proved}{HTML}{006100}
\definecolor{computed}{HTML}{1F4E78}
\definecolor{known}{HTML}{7F6000}
\definecolor{negative}{HTML}{9C0006}
\definecolor{open}{HTML}{7030A0}
\newcommand{\proved}{\textcolor{proved}{\textbf{[PROUVÉ DANS LA SOURCE]}}}
\newcommand{\computed}{\textcolor{computed}{\textbf{[CALCULÉ / CERTIFIÉ]}}}
\newcommand{\known}{\textcolor{known}{\textbf{[CONNU / À ATTRIBUER]}}}
\newcommand{\negative}{\textcolor{negative}{\textbf{[RÉSULTAT NÉGATIF]}}}
\newcommand{\openstatus}{\textcolor{open}{\textbf{[OUVERT / EXPLORATOIRE]}}}
\title{NvsNVP / Reverse-Hessian\\Résultats, certificats reproductibles, résultats négatifs et limites actuelles}
\author{Sasha Lempers}
\date{Release corrective v2.1.1 --- 19 juillet 2026}
\begin{document}
\maketitle

\begin{abstract}
Ce document organise le statut du programme Reverse-Hessian et de ses extensions calculatoires. Il relie deux manuscrits autonomes, cinq modules scientifiques, vingt-sept claims actifs, leurs certificats et leurs limitations. Il ne revendique ni une preuve de $VP\neq VNP$, ni une borne super-polynomiale de complexité déterminantale. La reproductibilité décrite ici vérifie des procédures, des fichiers et des calculs finis; elle ne remplace ni la revue experte des preuves, ni l'audit bibliographique de nouveauté.
\end{abstract}

\section{Convention de statut et non-claims}
\begin{itemize}[leftmargin=*]
  \item \proved{} : le paquet contient une démonstration textuelle; son contrôle mathématique indépendant reste nécessaire.
  \item \computed{} : un calcul exact ou modulaire est relié à un script et à un certificat.
  \item \known{} : l'énoncé est connu, reconstruit ou demande une attribution bibliographique précise.
  \item \negative{} : une famille d'invariants testée ne produit pas la séparation recherchée dans le régime indiqué.
  \item \openstatus{} : proposition, programme ou nouveauté non établie.
\end{itemize}

Les non-claims obligatoires sont : aucune preuve de $VP\neq VNP$; aucune borne super-polynomiale; aucune assimilation d'un test fini à une preuve générale; aucune assimilation d'un rang modulaire à une égalité sur $\mathbb{Q}$ sans certificat rationnel; aucune nouveauté déduite de la seule reproductibilité.

\section{Architecture scientifique}
Le papier \emph{A Uniform Line Formula and Hessian Orbit Incomparability} traite le coeur égal-taille. Le papier \emph{Reverse-Hessian VP/VNP Upgrade} traite le comportement sous padding et formule une barrière pour le déterminant de la Hessienne ordinaire. Ils restent autonomes parce qu'ils répondent à des questions différentes et portent des non-claims différents.

\section{Coeur égal-taille}
\proved{} Le claim \texttt{CORE-THM-001} renvoie à l'argument d'incomparabilité des clôtures d'orbites de $D(\det_n)$ et $D(\mathrm{per}_n)$, avec $D(f)=\det H_f$, pour le régime exact défini dans le manuscrit. La nouveauté n'est pas auditée par la présente release.

\computed{} Pour $n=3$, le certificat primaire donne les suites de rangs catalectiques
\[
(1,9,45,165,270,270,165,45,9,1)
\]
pour $D(\det_3)$ et
\[
(1,9,45,165,414,414,165,45,9,1)
\]
pour $D(\mathrm{per}_3)$. Ces valeurs sont régénérées et comparées byte-for-byte par le module core.

\section{Padding et barrière de la Hessienne ordinaire}
\negative{} L'extension VP/VNP explique pourquoi le déterminant de la Hessienne ordinaire se comporte mal sous le padding standard : l'invariant peut s'effondrer et ne fournit pas, à lui seul, une obstruction VP/VNP. Ce résultat de barrière est une conclusion scientifique positive sur la limite de la méthode, non une séparation de classes.

\section{Compound Hessian, Schur et Fitting}
\computed{} Le module compound vérifie des rangs catalectiques, des dimensions d'idéaux gradués, des noyaux avant minimalisation, des rangs de flattenings de Young et une borne modulaire pour $I_3$. Les certificats distinguent les égalités exactes, les bornes inférieures modulaires et les calculs non minimalisés.

\negative{} Dans les familles calculées, aucun test ne produit le saut de rang recherché en faveur de la permanente paddée; les valeurs permanentes sont plus petites dans les comparaisons enregistrées. La conclusion est strictement limitée aux familles et degrés testés.

\section{Renversement de rang paddé}
\computed{} Le module fournit un témoin exact $X_*$ avec $\mathrm{per}_3(X_*)=0$ et $\det H_{\mathrm{per}_3}(X_*)=-128$, les rangs de la Hessienne de $\det_4$ sur les strates de rang, et un certificat quotient--Fitting d'ordre neuf. Il conclut une non-inclusion pour le cas fixe documenté, sans revendiquer une conséquence asymptotique générale.

\section{Reconstruction asymptotique quotient--Fitting}
\proved{} Les sources donnent la borne de rang de la Hessienne sur $\det_n=0$, les formules le long de lignes et un critère de non-annulation pour le covariant d'ordre $2n+1$.

\known{} La borne $\lceil m^2/2\rceil$ de complexité déterminantale obtenue par cette reconstruction est connue dans la littérature. La release documente une reconstruction et son domaine d'efficacité, pas une nouvelle borne super-polynomiale.

\section{Programme jet-conormal}
\proved{} La première partie établit un rang $n^2$ pour un flattening de jet du déterminant dans le cadre énoncé.

\computed{} Des rangs fixes, des syzygies jacobiennes et des bornes modulaires conormales sont enregistrés par certificats.

\openstatus{} Le passage aux algèbres de Rees et au programme conormal supérieur reste exploratoire. Les rangs modulaires sont des bornes inférieures et ne constituent pas une séparation d'adhérences d'orbites.

\section{Reproductibilité}
Le profil \texttt{quick} vérifie le manifeste, la sécurité, les formats, les claims et les cinq modules. Le profil \texttt{full} ajoute les régénérations fraîches et les PDF déterministes. Le profil \texttt{heavy} tente les contrôles Sage lorsqu'ils sont disponibles. Toute dépendance absente reste explicitement marquée \texttt{BLOCKED\_NOT\_EXECUTED}.

\section{Limitations et questions ouvertes}
\begin{enumerate}[leftmargin=*]
  \item validation experte indépendante des démonstrations textuelles;
  \item audit bibliographique de la nouveauté des claims signalés;
  \item passage de cas fixes à des familles asymptotiques;
  \item remplacement des bornes modulaires par des certificats rationnels lorsque nécessaire;
  \item formalisation éventuelle dans un assistant de preuve;
  \item identification d'invariants stables sous padding capables de dépasser les barrières négatives observées.
\end{enumerate}

\section{Matrice des vingt-sept claims actifs}
\small
\begin{longtable}{@{}>{\raggedright\arraybackslash}p{1.6cm}>{\raggedright\arraybackslash}p{8.4cm}>{\raggedright\arraybackslash}p{2.5cm}>{\raggedright\arraybackslash}p{2.0cm}@{}}
\toprule
ID & Énoncé abrégé & Statut & Nouveauté \\
\midrule
\endfirsthead
\toprule ID & Énoncé abrégé & Statut & Nouveauté \\ \midrule
\endhead
CORE-THM-001 & Pour n>=3, le manuscrit canonique affirme l'incomparabilité des clôtures d'orbites de D(det\_n) et D(per\_n), où D(f)=det Hess(f). & calcul reproduit & non auditée \\
CORE-CERT-001 & Pour n=3, les rangs catalectiques certifiés de D(det\_3) et D(per\_3) sont respectivement (1,9,45,165,270,270,165,45,9,1) et (1,9,45,165,414,414,165,45,9,1). & calcul reproduit & non revendiquée \\
CORE-LIMIT-001 & Le noyau ne revendique ni obstruction standard pour le permanent paddé, ni borne super-polynomiale, ni VP!=VNP. & limite explicite & non revendiquée \\
COMP-001 & Pour det\_4 et z·per\_3, Phi\_2 survit au padding; les nombres de coordonnées non nulles sont 4032 et 540, avec 1152 et 213 générateurs distincts à scalaire près. & calcul reproduit & non revendiquée \\
COMP-002 & Les rangs catalectiques exacts d'ordres 0,1,2 sont [1,16,136] pour det\_4 et [1,10,55] pour z·per\_3. & calcul reproduit & non revendiquée \\
COMP-003 & Les dimensions exactes des pièces graduées de I\_2 en degrés 4,5,6 sont [625,5312,27608] pour det\_4 et [166,1908,12156] pour z·per\_3. & calcul reproduit & non revendiquée \\
COMP-SYZ-001 & Les dimensions des premières syzygies linéaires de I\_2 en degré 5, déduites exactement de la fonction de Hilbert, valent 4688 pour det\_4 et 748 pour z·per\_3. & calcul reproduit & non revendiquée \\
COMP-SYZ-002 & Les dimensions des noyaux quadratiques avant minimalisation en degré 6 valent 57392 pour det\_4 et 10420 pour z·per\_3. & calcul reproduit & non revendiquée \\
COMP-004 & Pour r=2, l'unique irréductible commun au domaine Sym\textasciicircum{}2(Sym\textasciicircum{}4 V) et à la cible composée est S\_(6,2)V; pour r=3 et r=4, les intersections exactes sont S\_(8,2,2)V et S\_(10,2,2,2)V, chacune de multiplicité un. & calcul reproduit & non revendiquée \\
COMP-005 & Les quatre flattenings de Young à deux cases donnent les rangs det/per: 136/55, 36/18, 136/55 et 120/45. & calcul reproduit & non revendiquée \\
COMP-006 & Pour I\_3 en degré 6, le rang exact côté z·per\_3 vaut 849 et le rang côté det\_4 est au moins 6832 sur Q, certifié par un rang modulo 1009. & calcul reproduit & non revendiquée \\
COMP-007 & Aucun des tests calculés dans cette branche ne fournit le saut de rang requis; tous les rangs certifiés sont plus faibles côté permanent paddé. & calcul reproduit & non revendiquée \\
RANK-001 & Au point entier X*=((-2,-1,-1),(-1,1,1),(-1,1,1)), per\_3(X*)=0 et det Hess(per\_3)(X*)=-128, donc le bloc hessien 9×9 de z·per\_3 est non nul sur son hypersurface. & calcul reproduit & non revendiquée \\
RANK-002 & Les rangs exacts du Hessien de det\_4 sur les strates de rang matriciel 0,1,2,3,4 sont 0,0,4,8,16; le rang maximal sur det\_4=0 est 8. & calcul reproduit & non revendiquée \\
RANK-003 & Pour le covariant quotient-Fitting d'ordre 9, le rang augmenté vaut N pour det\_4 et N+1 pour z·per\_3, avec N=10 150 394 132 736 000. & calcul reproduit & non revendiquée \\
RANK-004 & Il en résulte [z·per\_3]  not in  closure(GL\_16·[det\_4]). & preuve + certificat & non revendiquée \\
ASYM-001 & Sur l'hypersurface det\_n=0, le rang du Hessien est au plus 2n; les mineurs d'ordre 2n+1 sont donc divisibles par det\_n. & preuve incluse & non revendiquée \\
ASYM-002 & Pour X\_m(t)=J\_m+(t-1)E\_11, per\_m(X\_m(t))=(m-1)!(t+m-1) et det Hess(per\_m)(X\_m(t))=K\_m(t+m-3)\textasciicircum{}((m-2)\textasciicircum{}2), avec K\_m non nul. & borne modulaire & non revendiquée \\
ASYM-003 & Pour F\_\{n,m\}=ell\textasciicircum{}(n-m)per\_m, le covariant critique Theta\_\{2n+1\}(F\_\{n,m\}) est non nul si et seulement si 2n<m\textasciicircum{}2. & preuve incluse & non revendiquée \\
ASYM-004 & Si 2n<m\textasciicircum{}2, alors [ell\textasciicircum{}(n-m)per\_m]  not in  closure(GL\_\{n\textasciicircum{}2\}·[det\_n]). & preuve incluse & non revendiquée \\
ASYM-005 & La conséquence est la borne de complexité déterminantale bordée overline\{dc\}(per\_m) >= ceil(m\textasciicircum{}2/2). & preuve incluse & known result reconstructed \\
ASYM-006 & La méthode est sharp pour ce covariant: à 2n=m\textasciicircum{}2 le mineur actif maximal est divisible par la forme paddée, et pour 2n>m\textasciicircum{}2 tous les mineurs critiques s'annulent. & preuve incluse & non revendiquée \\
JET-001 & Pour n>=3 et 3<=t<=n, le flattening 1|(t-1) du t-ième jet de det\_n au point diag(I\_\{n-1\},0) a rang n\textasciicircum{}2. & preuve incluse & non auditée \\
JET-002 & Dans le cas det\_4 contre ell·per\_3, les rangs exacts des jets (ordre/split 2:1|1, 3:1|2, 4:1|3, 4:2|2) sont 8,16,16,36 contre 9,10,10,18. & calcul reproduit & non revendiquée \\
JET-003 & Les dimensions exactes des syzygies jacobiennes linéaires sont 30 pour det\_4 et 101 brutes pour ell·per\_3; après retrait des 96 sommants triviaux dus aux six générateurs nuls, il reste 5. & calcul reproduit & non revendiquée \\
JET-004 & Les premiers brins non saturés du graphe conormal ont été sondés modulo 1009, 10007 et 32003; les valeurs sont des bornes inférieures de rang, non des égalités rationnelles automatiques. & borne modulaire & non revendiquée \\
JET-005 & Les Hessiennes supérieures brutes saturent déjà côté déterminant dès l'ordre 3; la piste recommandée devient la résolution bigraduée de la Rees-algèbre jacobienne et les voisinages conormaux, décomposés par blocs de Schur. & proved first part exploratory second part & non revendiquée \\
\bottomrule
\end{longtable}

\normalsize

\section{Provenance et citation}
La source canonique est le tag \texttt{v2.1.1-corrective}. Le manifeste actif se trouve dans \texttt{07\_MANIFEST}. Les archives forensiques, l'audit initial et les anciens registres mixtes sont des assets de release avec SHA-256. Citer le manuscrit concerné, le module exact, le tag et, lorsqu'il existe, le DOI d'archive.

\begin{thebibliography}{9}
\bibitem{Valiant1979} L. G. Valiant, \emph{Completeness classes in algebra}, STOC, 1979.
\bibitem{BLMW2011} P. Bürgisser, J. M. Landsberg, L. Manivel, J. Weyman, \emph{An overview of mathematical issues arising in the Geometric Complexity Theory approach to VP vs. VNP}, SIAM J. Comput., 2011.
\bibitem{BIP2019} P. Bürgisser, C. Ikenmeyer, G. Panova, \emph{No occurrence obstructions in Geometric Complexity Theory}, JAMS, 2019.
\bibitem{Kumar2013} S. Kumar, \emph{Geometry of orbits of permanents and determinants}, Comment. Math. Helv., 2013.
\bibitem{Landsberg2017} J. M. Landsberg, \emph{Geometry and Complexity Theory}, Cambridge University Press, 2017.
\bibitem{MignonRessayre2004}
T. Mignon and N. Ressayre,
\emph{A quadratic bound for the determinant and permanent problem},
International Mathematics Research Notices, 2004.
\bibitem{LandsbergManivelRessayre2013}
J. M. Landsberg, L. Manivel and N. Ressayre,
\emph{Hypersurfaces with degenerate duals and the Geometric Complexity Theory Program},
Commentarii Mathematici Helvetici 88, 2013.
\end{thebibliography}
\end{document}
